La integraci\'on con resultados previos es el puente conceptual entre el an\'alisis global del repositorio y la priorizaci\'on intra-sistema. El proyecto ya construy\'o:

\begin{itemize}
\item nodos y regiones Mapper;
\item m\'etricas de sombra observacional por nodo;
\item TOI como prioridad regional;
\item ATI como prioridad de planetas ancla;
\item validaciones locales y multiescala sobre d\'eficit de vecinos.
\end{itemize}

En la capa regional, el proyecto usa el \'indice:
\[
TOI(v) =
Shadow(v)(1-I_{R^3}(v))C_{\mathrm{phys}}(v)S_{\mathrm{net}}(v).
\]
Este score combina frontera observacional, penalizaci\'on por dependencia de imputaci\'on en $R^3$, continuidad f\'isica y soporte de red. En la capa de ancla, el proyecto usa:
\[
ATI(p^\ast) =
TOI(v)\Delta_{\mathrm{rel,best}}(p^\ast)
(1-I_{R^3}(p^\ast))A(p^\ast).
\]
ATI hereda la prioridad regional y la modula con un d\'eficit local alrededor del planeta ancla, siempre con lectura prudente.

El nuevo m\'odulo no redefine estos \'indices. Los usa como se\~nal topol\'ogica exploratoria previa para evaluar un gap orbital entre dos planetas observados. Si el planeta interno o el externo pertenece a nodos con TOI alto, sombra alta o ATI alto, el intervalo hereda parte de esa prioridad. Una forma resumida de esta herencia es:
\[
T_{\mathrm{gap}} =
\max\{
TOI(v_{\mathrm{in}}),
TOI(v_{\mathrm{out}}),
ATI(p_{\mathrm{in}}),
ATI(p_{\mathrm{out}}),
Shadow(v_{\mathrm{in}}),
Shadow(v_{\mathrm{out}})
\}.
\]

La interpretaci\'on es directa: un gap grande con soporte topol\'ogico previo merece m\'as atenci\'on que un gap grande sin contexto adicional. Sin embargo, la relaci\'on sigue siendo indicativa y no causal. TOI, ATI y sombra no prueban que exista un planeta no detectado; solo refuerzan la idea de que la regi\'on local del cat\'alogo podr\'ia estar submuestreada.

Si estas tablas no existen, el reporte debe seguir siendo v\'alido. En ese caso, el an\'alisis cae al modelo orbital puro y la prioridad topol\'ogica del gap se fija en cero o queda marcada como no disponible. Este comportamiento es una decisi\'on metodol\'ogica deliberada para separar la evidencia orbital de la evidencia topol\'ogica previa.
